\documentclass[10pt,a4paper]{amsart}
\usepackage{amsmath,amssymb,amsthm,graphicx}
% Retained preamble for the Stacks Project chapter snapshot in this directory.
% The chapter's own preamble.tex is a copy of the shared Stacks Project preamble; the block
% below is sliced MECHANICALLY, line for line, from it. Nothing was added, removed or
% redefined; the wrapper class, the xr-hyper/hyperref wiring and the externaldocument list
% are omitted because the assembler supplies the document class.
%
% lines taken from preamble.tex: 3, 5, 9, 11, 12, 156, 157, 158, ... (51 lines total)
%
% package coverage restored from the chapter's own preamble (2026-10-13): xy, xy-2cell

\usepackage[all]{xy}
\xyoption{2cell}
\UseAllTwocells

\newenvironment{reference}{\comment}{\endcomment}
\newenvironment{slogan}{\comment}{\endcomment}
\newenvironment{history}{\comment}{\endcomment}
\theoremstyle{plain}
\newtheorem{theorem}[subsection]{Theorem}
\newtheorem{proposition}[subsection]{Proposition}
\newtheorem{lemma}[subsection]{Lemma}
\theoremstyle{definition}
\newtheorem{definition}[subsection]{Definition}
\newtheorem{example}[subsection]{Example}
\newtheorem{exercise}[subsection]{Exercise}
\newtheorem{situation}[subsection]{Situation}
\theoremstyle{remark}
\newtheorem{remark}[subsection]{Remark}
\newtheorem{remarks}[subsection]{Remarks}
\numberwithin{equation}{subsection}
\def\lim{\mathop{\mathrm{lim}}\nolimits}
\def\colim{\mathop{\mathrm{colim}}\nolimits}
\def\Spec{\mathop{\mathrm{Spec}}}
\def\Hom{\mathop{\mathrm{Hom}}\nolimits}
\def\Ext{\mathop{\mathrm{Ext}}\nolimits}
\def\SheafHom{\mathop{\mathcal{H}\!\mathit{om}}\nolimits}
\def\SheafExt{\mathop{\mathcal{E}\!\mathit{xt}}\nolimits}
\def\Sch{\mathit{Sch}}
\def\Mor{\mathop{\mathrm{Mor}}\nolimits}
\def\Ob{\mathop{\mathrm{Ob}}\nolimits}
\def\Sh{\mathop{\mathit{Sh}}\nolimits}
\def\NL{\mathop{N\!L}\nolimits}
\def\CH{\mathop{\mathrm{CH}}\nolimits}
\def\proetale{{pro\text{-}\acute{e}tale}}
\def\etale{{\acute{e}tale}}
\def\QCoh{\mathit{QCoh}}
\def\Ker{\mathop{\mathrm{Ker}}}
\def\Im{\mathop{\mathrm{Im}}}
\def\Coker{\mathop{\mathrm{Coker}}}
\def\Coim{\mathop{\mathrm{Coim}}}
\DeclareMathSymbol{\boxtimes}{\mathbin}{AMSa}{"02}
\def\QCohstack{\mathcal{QC}\!\mathit{oh}}
\def\Cohstack{\mathcal{C}\!\mathit{oh}}
\def\Spacesstack{\mathcal{S}\!\mathit{paces}}
\def\Quotfunctor{\mathrm{Quot}}
\def\Hilbfunctor{\mathrm{Hilb}}
\def\Curvesstack{\mathcal{C}\!\mathit{urves}}
\def\Polarizedstack{\mathcal{P}\!\mathit{olarized}}
\def\Complexesstack{\mathcal{C}\!\mathit{omplexes}}
\def\Pic{\mathop{\mathrm{Pic}}\nolimits}
\def\Picardstack{\mathcal{P}\!\mathit{ic}}
\def\Picardfunctor{\mathrm{Pic}}
\def\Deformationcategory{\mathcal{D}\!\mathit{ef}}
\usepackage{cleveref}
\crefname{originallemma}{original lemma}{original lemmas}
\Crefname{originallemma}{Original lemma}{Original lemmas}
\makeatletter
\newlabel{definition-quasi-coherent}{{definition-quasi-coherent}{source}{}{sourcepointer-1}{}}
\newlabel{definition-quasi-coherent@cref}{{[originallemma][0][]definition-quasi-coherent}{[1][0][]source}{}{}{}}
\makeatother
\crefname{originallemma}{original lemma}{original lemmas}
\Crefname{originallemma}{Original lemma}{Original lemmas}
\makeatletter
\newlabel{descent-definition-structure-sheaf}{{03DU}{source}{}{sourcepointer-2}{}}
\newlabel{descent-definition-structure-sheaf@cref}{{[originallemma][0][]03DU}{[1][0][]source}{}{}{}}
\makeatother
\crefname{originallemma}{original lemma}{original lemmas}
\Crefname{originallemma}{Original lemma}{Original lemmas}
\makeatletter
\newlabel{descent-lemma-sheaf-condition-holds}{{03DT}{source}{}{sourcepointer-3}{}}
\newlabel{descent-lemma-sheaf-condition-holds@cref}{{[originallemma][0][]03DT}{[1][0][]source}{}{}{}}
\makeatother
\crefname{originallemma}{original lemma}{original lemmas}
\Crefname{originallemma}{Original lemma}{Original lemmas}
\makeatletter
\newlabel{descent-proposition-equivalence-quasi-coherent}{{03DX}{source}{}{sourcepointer-4}{}}
\newlabel{descent-proposition-equivalence-quasi-coherent@cref}{{[originallemma][0][]03DX}{[1][0][]source}{}{}{}}
\makeatother
\crefname{originallemma}{original lemma}{original lemmas}
\Crefname{originallemma}{Original lemma}{Original lemmas}
\makeatletter
\newlabel{section-restriction}{{section-restriction}{source}{}{sourcepointer-5}{}}
\newlabel{section-restriction@cref}{{[originallemma][0][]section-restriction}{[1][0][]source}{}{}{}}
\makeatother
\crefname{originallemma}{original lemma}{original lemmas}
\Crefname{originallemma}{Original lemma}{Original lemmas}
\makeatletter
\newlabel{sites-modules-definition-site-local}{{03DL}{source}{}{sourcepointer-6}{}}
\newlabel{sites-modules-definition-site-local@cref}{{[originallemma][0][]03DL}{[1][0][]source}{}{}{}}
\makeatother
\title{Stacks, Sheaves on Algebraic Stacks: lemma-characterize-quasi-coherent-bis}
\author{The Stacks Project Authors}
\date{stacks-project:stacks-sheaves.tex:local-snapshot-20261012; GNU FDL 1.2 or later}
\begin{document}
\maketitle
\noindent\textit{Editorial note.} Statement span is lines 1818-1837 of sources/stacks-algebraic-stacks-sheaves/pinned-original-source.tex; proof span is lines 1839-1860, the complete author proof (20 body lines).\par\medskip
\noindent\textit{Editorial note.} Navigation stubs: 6 cited-but-not-included labels are supplied as non-load-bearing pointers with their pinned tags where the snapshot tag index has them.\par\medskip
\noindent\textit{Editorial note.} extra\_packages is empty on purpose: the retained preamble already loads xy with options and a second declaration clashes in the assembler wrapper.\par\medskip
\noindent\textit{Editorial note.} Text is retained verbatim, including the author's own typos and placeholders.\par\medskip
\noindent\textit{Editorial note.} Queue index 25 of queue\_algebraic-stacks-sheaves.json (pb=1148, stmt 1818-1837, proof 1839-1860).\par\medskip
\noindent\textit{Allowed standard prerequisites (editorial).}\begin{itemize}
\item Explicit prerequisite boundary. The cited inputs are prerequisites; the argument in the retained span is not replaced by any cited result.
\end{itemize}\medskip
\ifdefined\hypertarget\hypertarget{sourcepointer-1}{}\fi
\noindent\textit{Source reference.} Original lemma definition-quasi-coherent: sources/stacks-algebraic-stacks-sheaves/pinned-original-source.tex, source lines 1751--1760. Cited by the retained statement or proof but not part of the collected spans; supplied only as a navigation pointer. The argument does not depend on its content.\par\medskip
\ifdefined\hypertarget\hypertarget{sourcepointer-2}{}\fi
\noindent\textit{Source reference.} Original lemma 03DU: sources/stacks-algebraic-stacks-sheaves/pinned-original-source.tex, source lines 1860--1860. Cited by the retained statement or proof but not part of the collected spans; supplied only as a navigation pointer. The argument does not depend on its content.\par\medskip
\ifdefined\hypertarget\hypertarget{sourcepointer-3}{}\fi
\noindent\textit{Source reference.} Original lemma 03DT: sources/stacks-algebraic-stacks-sheaves/pinned-original-source.tex, source lines 1860--1860. Cited by the retained statement or proof but not part of the collected spans; supplied only as a navigation pointer. The argument does not depend on its content.\par\medskip
\ifdefined\hypertarget\hypertarget{sourcepointer-4}{}\fi
\noindent\textit{Source reference.} Original lemma 03DX: sources/stacks-algebraic-stacks-sheaves/pinned-original-source.tex, source lines 1860--1860. Cited by the retained statement or proof but not part of the collected spans; supplied only as a navigation pointer. The argument does not depend on its content.\par\medskip
\ifdefined\hypertarget\hypertarget{sourcepointer-5}{}\fi
\noindent\textit{Source reference.} Original lemma section-restriction: sources/stacks-algebraic-stacks-sheaves/pinned-original-source.tex, source lines 1185--1200. Cited by the retained statement or proof but not part of the collected spans; supplied only as a navigation pointer. The argument does not depend on its content.\par\medskip
\ifdefined\hypertarget\hypertarget{sourcepointer-6}{}\fi
\noindent\textit{Source reference.} Original lemma 03DL: sources/stacks-algebraic-stacks-sheaves/pinned-original-source.tex, source lines 1860--1860. Cited by the retained statement or proof but not part of the collected spans; supplied only as a navigation pointer. The argument does not depend on its content.\par\medskip
\begin{lemma}
\label{lemma-characterize-quasi-coherent-bis}
Let $p : \mathcal{X} \to (\Sch/S)_{fppf}$ be a category
fibred in groupoids. Let $\mathcal{F}$ be a presheaf of
modules on $\mathcal{X}$. The following are equivalent
\begin{enumerate}
\item $\mathcal{F}$ is an object of
$\textit{Mod}(\mathcal{X}_{Zar}, \mathcal{O}_\mathcal{X})$
and $\mathcal{F}$ is a quasi-coherent module on
$(\mathcal{X}_{Zar}, \mathcal{O}_\mathcal{X})$ in the sense of
Modules on Sites, Definition \ref{sites-modules-definition-site-local},
\item $\mathcal{F}$ is an object of
$\textit{Mod}(\mathcal{X}_\etale, \mathcal{O}_\mathcal{X})$
and $\mathcal{F}$ is a quasi-coherent module on
$(\mathcal{X}_\etale, \mathcal{O}_\mathcal{X})$ in the sense of
Modules on Sites, Definition \ref{sites-modules-definition-site-local}, and
\item $\mathcal{F}$ is a quasi-coherent module on $\mathcal{X}$
in the sense of Definition \ref{definition-quasi-coherent}.
\end{enumerate}
\end{lemma}

\begin{proof}
Assume either (1), (2), or (3) holds.
Let $x$ be an object of $\mathcal{X}$ lying over the scheme $U$.
Recall that $x^*\mathcal{F} = x^{-1}\mathcal{F}$ is just the
restriction to $\mathcal{X}/x = (\Sch/U)_\tau$ where
$\tau = fppf$, $\tau = \etale$, or $\tau = Zar$, see
Section \ref{section-restriction}.
By the definition of quasi-coherent modules on a ringed site
this restriction is quasi-coherent provided $\mathcal{F}$ is.
By Descent, Proposition \ref{descent-proposition-equivalence-quasi-coherent}
we see that $x^*\mathcal{F}$ is the sheaf associated to
a quasi-coherent $\mathcal{O}_U$-module and is therefore
a quasi-coherent module in the fppf, \'etale, and Zariski
topology; here we also use
Descent, Lemma \ref{descent-lemma-sheaf-condition-holds} and
Definition \ref{descent-definition-structure-sheaf}.
Since this holds for every object $x$ of $\mathcal{X}$,
we see that $\mathcal{F}$ is a sheaf in any of the three topologies.
Moreover, we find that $\mathcal{F}$ is quasi-coherent in any
of the three topologies directly from the definition of being
quasi-coherent and the fact that $x$ is an arbitrary object of $\mathcal{X}$.
\end{proof}

\end{document}
